% from ega0-19-fr.tex
\oldpage[0\textsubscript{IV-1}]{103}

\begin{corollary}[19.6.5.1]
\label{0.19.6.5.1-fr}
Soient $A$ un anneau local noethérien, $k$ un sous-corps de $A$ tel que $A$
soit une $k$-algèbre formellement lisse (pour sa topologie préadique). Alors
tout corps $K$ tel que $k\subset K\subset A$ est une extension séparable de
$k$.
\end{corollary}

\begin{proof}
En effet, pour toute extension finie $k'$ de $k$, l'anneau $K\otimes_k k'$ 
s'identifie à un sous-anneau de $A'=A\otimes_k k'$~; comme $A'$ est un anneau
régulier, il est réduit, donc il en est de même de $K$, ce qui prouve que $K$
est une extension séparable de $k$ (Bourbaki, \emph{Alg.}, chap.~VIII,
\textsection7, n\textsuperscript{o}~3, th.~1).
\end{proof}

On notera que cela démontre que la condition de l'énoncé de (19.6.1) est
\emph{nécessaire}.

\begin{remark}[19.6.5.2]
\label{0.19.6.5.2-fr}
Si $K$ est une extension \emph{non séparable} de $k$, l'anneau de séries
formelles $K[[T_1,\ldots,T_n]]$ \emph{n'est donc pas} une $k$-algèbre
formellement lisse (pour la structure de $k$-algèbre usuelle rappelée dans
(19.6.4)). D'autre part, il y a des $k$-algèbres formellement lisses $A$ qui
sont des anneaux locaux noethériens complets dont le corps résiduel $K$ est une
extension \emph{non séparable} de $k$ (par exemple les complétés des localisés
de $B=k[T_1,\ldots,T_n]$ en des idéaux maximaux $\mathfrak n$ tels que
$B/\mathfrak n$ soit une extension algébrique finie non séparable de $k$)~; un
tel anneau ne peut donc être $k$-isomorphe à $K[[T_1,\ldots,T_n]]$ bien qu'il
lui soit isomorphe en vertu de (19.6.5).

On peut dans certains cas préciser le corollaire (19.6.5)~:
\end{remark}

\begin{proposition}[19.6.6]
\label{0.19.6.6-fr}
Soient $k$ un corps, $A$ une $k$-algèbre locale noethérienne, $\mathfrak m$ son
idéal maximal, $A$ étant muni de la topologie $\mathfrak m$-préadique.
Supposons que le corps résiduel $K=A/\mathfrak m$ de $A$ soit tel qu'il existe
une extension radicielle finie $k_0$ de $k$ pour laquelle le corps résiduel de
l'anneau local $K\otimes_k k_0$ soit une extension séparable de $k_0$ (nous
exprimerons plus tard cette propriété en disant que $K$ est de
\emph{multiplicité radicielle finie} sur $k$ ($\mathbf{IV}$, 4.7.4) et nous
verrons que cette condition est toujours satisfaite lorsque $K$ est une
\emph{extension de type fini} de $k$). Alors les conditions suivantes sont
équivalentes~:
\begin{enumerate}
  \item[a)] $A$ est une $k$-algèbre formellement lisse.
  \item[b)] Il existe une extension radicielle finie $k'$ de $k$ telle que
  $\widehat A\otimes_k k'$ soit $k'$-isomorphe à une algèbre de séries
  formelles $K'[[T_1,\ldots,T_n]]$, où $K'$ est une extension séparable de
  $k'$.
  \item[b')] Il existe une extension finie $k'$ de $k$ telle que l'anneau
  semi-local $\widehat A\otimes_k k'$ ait au moins un anneau local composant
  qui soit $k'$-isomorphe à une algèbre de séries formelles
  $K'[[T_1,\ldots,T_n]]$, où $K'$ est une extension séparable de $k'$.
  \item[c)] $A$ est géométriquement régulier sur $k$ (19.6.5).
\end{enumerate}
\end{proposition}

\begin{proof}
Notons d'abord que si $k'$ est une extension radicielle de $k$, il n'y a qu'un
seul idéal de $A'=A\otimes_k k'$ au-dessus de $\mathfrak m$, formé des éléments
dont une puissance $p^h$-ième ($p$ exposant caractéristique de $k$) est dans
$\mathfrak m$ pour un $h$ convenable (Bourbaki, \emph{Alg. comm.}, chap.~V,
\textsection2, n\textsuperscript{o}~3, lemme~4)~; $A'$ est donc un anneau
local, et il en est de même de
$K\otimes_k k'=(A\otimes_k k')/(\mathfrak m\otimes_k k')$~; en outre les corps
résiduels de ces deux anneaux sont identiques. Rappelons d'autre part que si
$K$ est une extension séparable de $k$, alors, pour toute extension finie
$k''$ de $k$, $K\otimes_k k''$ est composé direct de corps (Bourbaki,
\emph{Alg.}, chap.~VIII, \textsection7, n\textsuperscript{o}~3, cor.~1 du
th.~1), et par suite $\mathfrak m\otimes_k k''$ est le radical de
$A\otimes_k k''$, et les corps composants de $K\otimes_k k''$ sont les corps
résiduels aux idéaux maximaux de $A\otimes_k k''$~; en

